\documentclass{article}
\usepackage{amsmath}
\usepackage{amssymb}
\usepackage{array}
\usepackage{booktabs}
\usepackage{textcomp}
\usepackage{graphicx}

\usepackage{amsthm}
\newtheorem{theorem}{Theorem}[section]
\usepackage{algpseudocodex}

\newcommand{\sol}{\textbf{sol.}}

\newtheorem{lemma}[theorem]{Lemma}



\newcommand{\coNP}{\textrm{co-}\mathbb{NP}}

\title{sesh-meeting-2026-02-04}
\author{Ian Chi}
\date{\today}
\begin{document}
\maketitle

given n coins, each with probability \(p\) of flipping heads, what is the probability that all of them flip heads. this is equal to \((p)^n\). 

(in this case, \(p\) is analogous to all \(x_{n}\) expected values being under 2/3, sampling each \(x_n\) \(m\) times, given \(m\) is sufficiently large, what is the probability that one sample is above 2/3. this is the \(a\) light detection case. \(p\) is thus close to 1.)


if i want to use union bound, the union of all events being 1 is less than the sum. each sample goes from \(p\) probability to \(\frac{n+1}{n}p\). it requires \(n\) extra samples by chernoff? since \(2/3\) is fixed and \(p\), the expected value is known, we can use multiplicative percentage error





I want \(\left( p \right)^n >2/3\). 1 extra sample requires \(p\) to increase by \(\alpha\). to find \(\alpha\), 

$$\Delta p = \left( \frac{2}{3} \right)^{\frac{1}{n+1}} - \left( \frac{2}{3} \right)^{\frac{1}{n}} = -\frac{2}{3}^{1/n}$$

\begin{equation}
    \begin{split}
        \frac{2}{3} &= \left( 1-p \right)^n \\
    \end{split}
\end{equation}



\end{document}